% Created 2026-09-05 Sat 13:25
% Intended LaTeX compiler: xelatex
\documentclass[12pt]{article}
\usepackage[utf8]{inputenc}
\usepackage[T1]{fontenc}
\usepackage{graphicx}
\usepackage{longtable}
\usepackage{wrapfig}
\usepackage[margin=1in]{geometry}
\usepackage{rotating}
\usepackage{setspace}
\onehalfspacing
\usepackage[normalem]{ulem}
\usepackage{amsmath}
\usepackage{parskip}
\usepackage{amssymb}
\usepackage{capt-of}
\setcounter{secnumdepth}{4} % Tells LaTeX to number down to paragraphs
\setcounter{tocdepth}{4}
\usepackage{hyperref}
\usepackage[style=oscola, backend=biber]{biblatex}
\addbibresource{/Users/krishnajani/Documents/Libib.bib}
\usepackage{csquotes}
\usepackage[british]{babel}
\setlength{\parskip}{0.5\baselineskip}
\author{Krishna Jani}
\date{\today}
\title{Financial Lease - HP v. Nufuture}

\begin{document}

\maketitle

S. 5(9) - A financial debt is a debt, alongside interest \emph{if any} which is disbursed against consideration for \emph{time-value of money.}
\section{Essentials of Financial Debt}
\label{sec:org82e7169}

\noindent The various forms in which a debt can manifest is provided under S. 5(8) as follows
\begin{enumerate}
\item Borrow against interest
\item Borrow under a credit facility
\item Borrow under a note purchase, loan stock, debentures, bonds instrument
\item Lease or Hire Contract for \emph{financial} or \emph{capital lease} under IND AS
\item Receivables sold or discounted.
\item Derivatives
\item Counter-indemnity - These are cases where the guarantor will also enter into a contract with the debtor to ensure that it is able to recover the loss it undertakes as a guarantor
\item Liability under a contract of guarantee or indemnity
\end{enumerate}
\subsubsection{Financial Lease}
\label{sec:orge839971}
A lease allows the lessee to use and enjoy the underlying asset on the payment of a lease amount. This lease can be both \textbf{financial} and \textbf{operational} in nature. 
Where it is operational, the lessee only has the rights to \uline{possess and use the underlying asset}. Where it is a financial lease, \uline{risk and rewards associated with ownership} may also be transferred.
The objective of a financial lease is to eventually over a period of time acquire the property and is often used for high-value assets. Almost, like for instance, a bareboat charter with a purchase option (Poompuhar Shipping Company)

\noindent \textbf{Requirements for a FL (generally})
\begin{enumerate}
\item Lessee will transfer the ownership to the lessee at the end of the lease period
\item Lessee has the option to purchase the underlying asset at a price sufficiently lower than its FMV
\item The lease term is for a major part of the economic lifetime of the underlying asset
\item At inception the value of lease payments would amount to the substantial value of the underlying asset. (IFRS 16)
\end{enumerate}
\subsubsection{HP v. Nufuture}
FC filed a S.7 application for initiating CIRP. Both parties had entered into a ``MFRA'' Agreement which was for leasing IT Equipment. There was a default in the payment terms which triggered the S. 7 application.

HP contended that it was a financial lease, since as per IndAS 116 a financial lease is one where substantial risks and rewards associated with the asset are transferred. It also submitted that where the value of the lease payment covers substantially all of the value of the asset, it constitutes a financial lease. On this ground, it was submitted that the lease amount covers 90\% of the asset cost and is thus \textbf{financial lease}.

Nufuture argued that risks associated with owership were not transferred. It did not have the option to purchase. The Agreement stipulated return of equipment. In case of default, HP could seek possession of the leased equipment. Thus it constitutes an \textbf{operational lease}

\paragraph{NCLT Order} --- It was held that the lease is a \textbf{financial lease} and thus a Section 7 application is valid. Since the lease value is 90\% of the asset value it fulfills IndAS requirements. Merely because the ownership is not transferred or the right to re-possession exists, does not derogate from its nature as a financial debt.

\subsubsection{Evidence of Default}

Section 7(3) requires the FC to furnish the \textbf{evidence of default} with IU or Regulators. IU is not the only way to do so. Even \textbf{books of account, tribunal order evidencing debt} can be filed with the AA.

\paragraph{Vijay Kumar Singhania v. BoB} --- It clarified the evidentiary requkrements under S. 7

\begin{enumerate}
  \item While filing with IU is key, it is not mandatory. FC can also provide other evidence. 
  \item The rules do not override S. 7 in so far as the FC is allowed to also submit alternative evidence
  \item Existance of counterclaim does not disentitle the FC to initiate CIRP. 
\end{enumerate}
The position has been finally upheld by the Supreme  Court in Swiss Ribbons Case.\footcite{justicer.fnarimanSwissRibbonsPvt2019}
\end{document}
